\corrigesChapitre{1}{Matière : mesurer, séparer, identifier}
\par\noindent\begin{minipage}{\linewidth}
\begin{corrige}{1}{Trois états}
\begin{questions}
\item Le solide, hors déformation imposée.
\item Un gaz se répartit dans tout le volume disponible.
\end{questions}
\end{corrige}
\end{minipage}
\par\noindent\begin{minipage}{\linewidth}
\begin{corrige}{2}{Une eau salée}
\begin{questions}
\item Mélange : eau et espèces dissoutes.
\item Non ; plusieurs espèces peuvent former une seule phase visible.
\end{questions}
\end{corrige}
\end{minipage}
\par\noindent\begin{minipage}{\linewidth}
\begin{corrige}{3}{Les bons mots}
\begin{questions}
\item Vaporisation.
\item Liquéfaction ; on rencontre aussi condensation liquide.
\end{questions}
\end{corrige}
\end{minipage}
\par\noindent\begin{minipage}{\linewidth}
\begin{corrige}{4}{Volumes}
\begin{questions}
\item $350\unite{mL}$.
\item $18\unite{mL}$.
\end{questions}
\end{corrige}
\end{minipage}
\par\noindent\begin{minipage}{\linewidth}
\begin{corrige}{5}{Le quotient}
\begin{questions}
\item $40/50=0{,}80\unite{g/mL}$.
\item $800\unite{kg/m^3}$.
\end{questions}
\end{corrige}
\end{minipage}
\par\noindent\begin{minipage}{\linewidth}
\begin{corrige}{6}{Solide immergé}
\begin{questions}
\item $V=43-28=15\unite{cm^3}$.
\item $\rho=40{,}5/15=2{,}70\unite{g/cm^3}$.
\end{questions}
\par\smallskip\noindent\textbf{Pour comprendre.}\enspace $V=43-28=15\unite{mL}=15\unite{cm^3}$. On calcule ensuite $\rho=40{,}5/15=2{,}70\unite{g/cm^3}$. Utiliser 43 mL compterait aussi l’eau initiale.\par
\pcFigureEx{figures/rev-cor-immersion.pdf}{Le solide déplace 15 mL ; on utilise ce volume avec la seule masse du solide.}
\end{corrige}
\end{minipage}
\par\noindent\begin{minipage}{\linewidth}
\begin{corrige}{7}{Une masse prévue}
\begin{questions}
\item $m=0{,}90\times250=225\unite g$.
\item Elle diminue puisque $m/V$ diminue.
\end{questions}
\end{corrige}
\end{minipage}
\par\noindent\begin{minipage}{\linewidth}
\begin{corrige}{8}{Retrouver un volume}
\begin{questions}
\item $V=178/8{,}9=20\unite{cm^3}$.
\item Le volume extérieur comprend la cavité ; ce n’est pas le seul volume de cuivre.
\end{questions}
\par\smallskip\noindent\textbf{Pour comprendre.}\enspace Partir de $m=\rho V$ et isoler $V=m/\rho=178/8{,}9=20\unite{cm^3}$. Le volume extérieur d’un objet creux ne se substitue pas au volume du matériau.\par
\pcFigureEx{figures/cor-coque.pdf}{Une cavité contribue au volume extérieur, mais pas au volume de cuivre.}
\end{corrige}
\end{minipage}
\par\noindent\begin{minipage}{\linewidth}
\begin{corrige}{9}{Une limite de dissolution}
\begin{questions}
\item Environ 36 g se dissolvent et 4 g restent solides à l’équilibre, dans ce modèle.
\item $36\times2{,}5=90\unite g$.
\end{questions}
\end{corrige}
\end{minipage}
\par\noindent\begin{minipage}{\linewidth}
\begin{corrige}{10}{Interpréter un palier}
\begin{questions}
\item L’énergie transférée sert notamment au changement d’état pendant la coexistence liquide-vapeur.
\item Non : de la vapeur quitte le système pesé ; inclure cette matière dans le bilan.
\end{questions}
\pcFigureEx{figures/rev-cor-ebullition.pdf}{La température reste constante pendant le changement d’état, à la pression annoncée.}
\end{corrige}
\end{minipage}
\par\noindent\begin{minipage}{\linewidth}
\begin{corrige}{11}{Choisir une séparation}
\begin{questions}
\item Filtrer le sable, puis distiller le filtrat avec un dispositif adapté.
\item Non : la potabilité nécessite des critères et des analyses supplémentaires.
\end{questions}
\pcFigureEx{figures/rev-cor-separation.pdf}{Filtration pour retirer le sable ; distillation du filtrat pour recueillir l’eau et laisser le sel non volatil.}
\end{corrige}
\end{minipage}
\par\noindent\begin{minipage}{\linewidth}
\begin{corrige}{12}{Une mesure discutable}
\begin{questions}
\item $54/25=2{,}16\unite{g/mL}$, trop faible.
\item Oui : la bulle peut augmenter le volume déplacé apparent.
\end{questions}
\par\smallskip\noindent\textbf{Pour comprendre.}\enspace Avec le volume apparent, $54/25=2{,}16\unite{g/cm^3}$. Avec le volume réel de 20 cm³, on aurait $54/20=2{,}70\unite{g/cm^3}$. Une mesure biaisée peut donc changer l’identification proposée.\par
\pcFigureEx{figures/cor-bulle.pdf}{Une bulle fait attribuer au solide un volume trop grand et une masse volumique trop petite.}
\end{corrige}
\end{minipage}
\par\noindent\begin{minipage}{\linewidth}
\begin{corrige}{13}{Identifier avec une plage}
\begin{questions}
\item La valeur 2,70 appartient à l’intervalle 2,63 à 2,73 : résultat compatible.
\item Non ; d’autres matériaux ou compositions peuvent convenir. Un second test serait utile.
\end{questions}
\pcFigureEx{figures/rev-cor-intervalle.pdf}{2,70 appartient à l’intervalle 2,63 à 2,73 g/cm³ ; compatibilité ne signifie pas identification certaine.}
\end{corrige}
\end{minipage}
\par\noindent\begin{minipage}{\linewidth}
\begin{corrige}{14}{Glace et océans}
\begin{questions}
\item Cette eau peut rejoindre l’océan depuis un réservoir auparavant sur les continents.
\item La dilatation thermique de l’eau de mer augmente son volume à masse donnée.
\end{questions}
\pcFigureEx{figures/cor-glaces.pdf}{Distinguer apport d’eau par les glaciers continentaux et dilatation de l’eau ; la banquise flottait déjà.}
\end{corrige}
\end{minipage}
\par\noindent\begin{minipage}{\linewidth}
\begin{corrige}{15}{Construire une preuve}
\begin{questions}
\item Mesurer leurs masses pour un volume connu et comparer leurs masses volumiques ; utiliser les protections prescrites.
\item Une différence de température ou de méthode peut produire un écart qui ne vient pas de la composition.
\end{questions}
\end{corrige}
\end{minipage}
\par\noindent\begin{minipage}{\linewidth}
\phantomsection\label{cor-defi-1}\pcSolutionDefi{1}{Le flacon inconnu}\begin{questions}
\item $m=109-45=64\unite g$ ; $\rho=64/80=0{,}80\unite{g/mL}$.
\item L’éthanol est le plus proche, sans preuve de pureté.
\item La température ou une lecture grossière peut expliquer l’écart. Répéter avec verrerie adaptée et comparer plusieurs propriétés.
\end{questions}
\par\smallskip\noindent\textbf{Pour comprendre.}\enspace La masse du liquide vaut $109-45=64\unite g$. Avec 80 mL, $\rho=64/80=0{,}80\unite{g/mL}$. La proximité avec 0,79 g/mL rend l’éthanol compatible ; d’autres liquides ou mélanges peuvent avoir une valeur proche.\par
\pcFigureEx{figures/rev-cor-tare.pdf}{Soustraire la masse du flacon avant de former le quotient masse/volume.}
\end{minipage}
\par\noindent\begin{minipage}{\linewidth}
\phantomsection\label{cor-activite-1}\pcSousTitre{Pistes pour l’investigation}
Une tare correcte soustrait la masse du récipient. Pour un liquide homogène à température stable, m est sensiblement proportionnelle à V ; les écarts doivent être comparés à la lecture. Pour un solide insoluble, mesurer la masse puis le volume déplacé avec immersion complète et sans bulle. Une identification par masse volumique est compatible avec une référence, sans prouver l’unicité ni la pureté.\par\end{minipage}
